% Generated from 07-conclusion-and-future-work.md; edit Markdown, then rebuild.
\chapter{总结与展望}\label{chap:7}

本文围绕具有固定验收入口的 Python 代码修复任务，在已有终端编程智能体基础上，设计并实现了多智能体协作与验证系统。研究重点是使文件修改、命令执行、角色交接与任务结果具有可检查的关联，并通过功能测试和真实修复运行评价其行为。本章总结已完成的工作，回应绪论提出的研究问题，并说明当前局限与后续研究方向。

\section{工作总结}\label{sec:7.1}

首先，系统保留了编程智能体的模型交互、工具调用和代码定位能力，并在其上建立专用修复流程。协调、修复和验证角色使用独立上下文，按顺序完成任务分析、源码修改与候选检查。工具注册表根据职责收窄能力，只有修复角色可以修改业务源码；CLI、REPL、SDK 和 Runtime 通过各自的适配方式接入同一修复编排器，形成可实际调用的系统。

其次，系统建立了原始源码、修复工作副本、验收副本和交付目录之间的边界。初始文件快照和来源指纹确定候选参照，文件比较识别新增、修改及删除，并拒绝支持范围内的测试和关键配置变更。模型命令由受限 Docker 容器执行，固定验收在重建的候选副本中运行。成功候选以补丁交付，系统不自动覆盖原仓库；超时、取消和异常结束均有相应的资源回收处理。

再次，系统将基线结果、候选差异、真实测试、评审意见与终止状态组织到同一次运行中。固定验收是程序执行的判断依据，评审文字提供解释和后续建议。候选失败时可以在剩余预算内继续尝试；调用次数、token 阈值、尝试次数和时限共同控制流程，无进展、环境错误和预算耗尽分别记录。补丁、报告和事件文件为成功与未完成结果提供检查入口。

同时，系统接入 Skill 与项目 Memory。角色技能描述复现、最小修改和验证步骤，报告保留技能版本与内容摘要；项目记忆按作用域召回，对修复经验增加来源指纹和有效期约束，仅在候选通过验收且流程成功时自动保存。两类知识作为角色输入，不能改变固定验收或直接覆盖运行状态。

最后，本文完成了工程验证和受控修复实验。归档的 160 项自动化测试包含既有回归及新增修复测试，覆盖副本保护、固定验收、有限重试、容器清理和接口适配。实验基于固定 python-dotenv 源码构造 10 个人工缺陷，将其中 2 个用于开发、8 个用于保留评测，完成三种策略共 24 次真实模型运行，并使用运行后加入的额外测试检查已导出候选。

\section{主要结论}\label{sec:7.2}

\subsection{通用编程能力可以与专用执行边界结合}\label{sec:7.2.1}

针对如何保留编程核心并控制修复影响范围的问题，本文采用复用查询循环、重新组织工具注册及独立管理副本的方式。模型仍能读取、检索、编辑和执行辅助命令，而修改对象与命令可见范围由执行层确定。功能测试和实际交付记录支持该方案在当前 Python 任务范围内可运行，原始源码不因自动修复直接改变。

执行边界需要分层落实：宿主文件工具检查路径与文件类型，候选交接检查修改范围，命令执行限制容器资源和挂载。容器只是其中一个组成部分，不能代替文件工具约束，也不能据此推断所有通用入口均已获得相同隔离。本文验证的是所实现流程及其检查场景，不构成对任意恶意代码的完整安全证明。

\subsection{执行证据与流程状态需要分别保留}\label{sec:7.2.2}

针对如何组织角色交接和判断完成的问题，本文以实际候选及固定命令结果连接各角色。验证角色的肯定意见无法覆盖真实失败；基线已通过、环境不可用或没有有效修改时，也不会统一输出修复成功。这使调用者能够区分模型陈述、文件变化和可执行检查。

候选通过测试与流程成功完成同样需要区分。实验中有一次验证反馈运行的候选通过固定及额外测试，但验证角色调用期间预算耗尽，最终仍为失败。保留这两个结果，可以说明候选满足了哪些条件，以及策略在哪个阶段未完成。只汇总模型文字或只检查最终补丁，都会遗漏一部分流程成本和终止信息。

\subsection{多角色收益取决于任务与预算条件}\label{sec:7.2.3}

针对相同预算下不同策略的表现问题，单角色执行反馈、评审反馈多角色和验证反馈多角色的流程成功数分别为 8、5、4，额外检查通过数分别为 8、5、5。各组未观察到错误验收，两个多角色组共 7 次失败均由 token 阈值触发。多角色平均累计 token 约为单角色的 2.07 倍和 1.93 倍。

这些结果没有显示多角色在本任务集上的成功率优势，反映出局部缺陷中协调、重复阅读和复查可能先消耗可用预算。三组均保留固定测试，零错误验收也不能单独归因于验证角色。所有运行只建立了一轮尝试，因而尚不能从真实模型实验判断多轮反馈的收益。

本文的主要成果体现为可运行的修复架构、明确的执行约束和可核对的结果组织，以及对协作成本的实际记录。该结论限定于当前小规模任务与设置，不推广为多智能体普遍无效，也不把功能完整直接等同于修复质量领先。

\section{当前局限}\label{sec:7.3}

任务方面，正式评测仅有同一项目的 8 个人工缺陷，集中于少量文件，任务描述提供了较明确的行为线索。测试由本项目编写，额外检查也只有有限覆盖，未经过独立专家确认。现有结果不足以代表跨文件复杂问题、真实历史缺陷或未知项目的整体表现。

实验方面，每个任务和策略只运行一次，条件顺序未随机化，远程模型服务与同机负载也未被完全控制。统一 token 阈值对不同角色数量的策略形成不同约束，结果反映当前预算下的整体流程表现，不能直接推断调整预算后的排序。由于没有进入第二轮，实验对反馈重试效果的评价尤其有限。

知识方面，真实实验始终开启 Skill、关闭跨任务 Memory，没有单独比较技能选择、经验召回或记忆更新对质量的影响。已有功能测试支持它们能够按设计接入，但不支持独立增益结论。源码指纹也只覆盖支持范围内的文件，不能完整识别镜像、依赖和运行环境变化。

工程方面，当前采用顺序协作和单个写入角色，未实现并行候选生成与补丁冲突合并。快照和补丁处理受文件类型、大小与固定测试规则限制，容器依赖预先准备的环境。报告支持结果检查和过程分析，但没有完整保存模型原始交互，也不具备任意进程崩溃后的自动恢复能力。

\section{后续研究方向}\label{sec:7.4}

\subsection{扩展独立任务与评估设计}\label{sec:7.4.1}

后续可优先建立包含更多仓库、跨文件依赖和真实缺陷的冻结任务集，并由独立人员检查任务及验收条件。开发与评测应继续分离，在运行前明确预算、成功定义和额外检查，再采用重复运行与条件顺序随机化报告结果分布。对成本、成功和未检查样本保持统一统计口径，能够减少只选成功案例造成的偏差。

更复杂任务还应包含首轮候选真实失败、但存在合理修复空间的场景，用于检验模型是否能够利用后续证据。脚本化测试继续负责稳定检查流程分支，真实模型实验则负责评估实际改进，两类证据共同使用但分别解释。

\subsection{研究按风险选择角色与预算分配}\label{sec:7.4.2}

当前实验提示，并非每个局部修改都需要同等规模的规划和复查。后续可研究依据变更范围、测试覆盖和失败类型选择角色，并在任务开始前规定角色预算与保留额度，减少前置阶段耗尽资源而无法完成交接的情况。上下文摘要和按需检索也可作为降低重复输入的方向。

这些方案需要在新的冻结设置下比较完成率、错误验收和累计用量，不能仅通过减少调用就认定有效。若改变测试通过后是否必须继续评审，还应同步重定义流程完成条件，并与原策略分别报告，避免事后改变成功口径。并行候选可作为进一步研究，但需要独立工作副本、候选归属及冲突处理机制，不宜直接在共享目录中增加多个写入角色。

\subsection{独立评价技能与记忆的作用}\label{sec:7.4.3}

后续可分别设置技能开启与关闭、记忆开启与关闭的对照条件，控制其他变量，检查知识输入是否改善修复，或者带来过时建议与额外成本。记忆实验还应按项目和时间划分任务，避免把同一缺陷答案泄漏到评测输入，并同时记录无关召回、错误经验和失效过滤。

经验适用条件可进一步纳入环境与依赖指纹，结合来源报告重新验证。技能选择也可从固定角色映射扩展为可检查的任务匹配机制。这些都属于尚待验证的研究方向，本轮功能接入不预先保证其质量收益。

\subsection{完善执行记录与交付核对}\label{sec:7.4.4}

后续可以在明确保留范围和敏感信息处理规则的前提下，补充环境元数据、模型请求摘要和更细的阶段用量，增强失败定位能力。对报告写入中断、补丁重放和人工应用前的来源检查，也可以建立更明确的交付协议。记录的目标应是支持核对，而不是无边界保存全部源码与凭据。

本文形成了从编程智能体能力到受约束修复任务的一条实现路径，并通过实际运行保留其成功和失败。进一步研究应继续以可执行检查和可复核记录评价方案，使角色数量、知识输入与资源投入能对应到明确的任务收益。
